\section{Central-factorial kernels and finite zeta reductions}
\label{rlc:sec:central}

\subsection{The logistic convolution polynomial}
Define monic polynomials
\begin{align}
 p_{2m}(X)&=X\prod_{j=1}^{m-1}\bigl(X^2+4j^2\pi^2\bigr),
          &&m\ge1,\label{rlc:eq:peven}\\
 p_{2m+1}(X)&=\prod_{j=1}^{m}\bigl(X^2+(2j-1)^2\pi^2\bigr),
          &&m\ge0.\label{rlc:eq:podd}
\end{align}
Thus $p_1=1$, $p_2=X$, and
\begin{equation}\label{rlc:eq:precurrence}
 p_{r+2}(X)=(X^2+r^2\pi^2)p_r(X).
\end{equation}
These are central-factorial product polynomials after an imaginary
rescaling of their variable. Their degree is $r-1$ and their parity
is $(-1)^{r-1}$.

\begin{lemma}[All-depth elementary kernel]\label{rlc:lem:logistic}
For every $r\ge1$ and real $\mu$,
\begin{equation}\label{rlc:eq:logistic}
 f_0^{*r}(\mu)=
 \frac{p_r(\mu)}{(r-1)!\,[1-(-1)^r e^{-\mu}]}.
\end{equation}
At even depth the quotient at $\mu=0$ is filled in by continuity.
\end{lemma}
\begin{proof}
At depth one this is $f_0$. At depth two the substitution $y=e^x$
gives $f_0*f_0(\mu)=\mu/(1-e^{-\mu})$, including its limit $1$.
Let $g(p)=\pi\csc(\pi p)$. Direct differentiation gives
\[
 \frac{\dd^2}{\dd p^2}g(p)^r
  =r(r+1)g(p)^{r+2}-r^2\pi^2g(p)^r.
\]
Multiplication by $\mu^2$ corresponds to two derivatives of the
bilateral transform. Transform uniqueness therefore implies
\[
 f_0^{*(r+2)}(\mu)=
 \frac{\mu^2+r^2\pi^2}{r(r+1)}f_0^{*r}(\mu).
\]
The two initial cases and \eqref{rlc:eq:precurrence} prove the result.
All transforms used here lie in the common strip $0<\Rea p<1$;
polynomial multiplication is justified by the neighboring-strip
estimates from Lemma~\ref{rlc:lem:weighted}.
\end{proof}

\subsection{A finite formula away from the central point}
Put $\epsilon_r=(-1)^r$.
\begin{theorem}[Finite Lerch closure]\label{rlc:thm:lerch}
For $\mu<0$, $\Rea a>0$, and $W\in\C$,
\begin{equation}\label{rlc:eq:lerchclosure}
 \boxed{\displaystyle
 \CC_{r,a}(W;\mu)=
 \frac{-\epsilon_r e^\mu}{(r-1)!}
 \sum_{j=0}^{r-1}\frac{(-1)^j p_r^{(j)}(\mu)}{j!}
 (W)_j\Phi(\epsilon_r e^\mu,W+j,a).}
\end{equation}
The combined expression continues to all real $\mu$. At even depth and
$\mu>0$, the continuation of the \emph{whole sum} is independent of
whether the Lerch cut is approached from above or below.
\end{theorem}
\begin{proof}
For $\Rea W>0$, insert \eqref{rlc:eq:logistic} into
\eqref{rlc:eq:basefractional}. When $\mu<0$, expand the denominator as
\[
 \frac1{1-\epsilon_r e^{u-\mu}}
 =-\epsilon_r e^{\mu-u}
       \sum_{n=0}^\infty\epsilon_r^n e^{n(\mu-u)}.
\]
Expand $p_r(\mu-u)$ by its finite Taylor formula. The remaining
integrals are
$\Gamma(W+j)(n+a)^{-W-j}/\Gamma(W)$, yielding
\eqref{rlc:eq:lerchclosure}. Absolute convergence justifies the steps.
Both sides are entire in $W$ for fixed $\mu<0$.

The left side is holomorphic for $|\Ima\mu|<r\pi$. At even depth,
the discontinuity of the Lerch integral at $z=e^\mu>1$ is, up to a
common nonzero sign convention,
$2\pi i e^{-a\mu}\mu^{W+j-1}/\Gamma(W+j)$.
After multiplication by $(W)_j$ and the other factors, the jump of the
sum is proportional to
\[
 \sum_{j=0}^{r-1}\frac{(-1)^j p_r^{(j)}(\mu)\mu^j}{j!}
 =p_r(0)=0.
\]
This calculation is initially made away from Gamma poles and then
continued in $W$. Thus the two real boundary values agree and coincide
with the ordinary convolution. At $\mu=0$ use the entire convolution,
or the central formula below. At odd depth the argument
$-e^\mu$ does not cross the positive Lerch cut.
\end{proof}

At depth two and $a=1$, the formula becomes
\begin{equation}\label{rlc:eq:fugacity2}
 \int_0^\infty\Li_s(-y)\Li_t(-e^\mu/y)\frac{\dd y}{y}
 =W\Li_{W+1}(e^\mu)-\mu\Li_W(e^\mu).
\end{equation}
This is a branch-cancelled combination, not a license to discard the
imaginary parts of its terms separately. Its limit at $\mu=0$ is
$\E_1(W)$.

\subsection{Pole-free central formulas}
Define the entire alternating Hurwitz function
\begin{equation}\label{rlc:eq:eta}
 \eta(v,a)=\Phi(-1,v,a)
 =2^{-v}\left[\zeta\left(v,\frac a2\right)
                -\zeta\left(v,\frac{a+1}2\right)\right].
\end{equation}
The two residues at $v=1$ cancel. In particular
$\eta(0,a)=1/2$ and
$\eta(1,a)=\tfrac12[\psi((a+1)/2)-\psi(a/2)]$.
Write
\[
 p_{2m}(X)=\sum_{k=0}^{m-1}A_{m,k}X^{2k+1},\qquad
 p_{2m+1}(X)=\sum_{k=0}^{m}B_{m,k}X^{2k}.
\]
The coefficients are explicitly the elementary symmetric functions
of the squared even or odd multiples of $\pi$ in
\eqref{rlc:eq:peven}--\eqref{rlc:eq:podd}.

\begin{theorem}[Central finite closure]\label{rlc:thm:central}
For $\Rea a>0$ and every $W\in\C$,
\begin{align}
 \CC_{2m,a}(W;0)
 &=\frac1{(2m-1)!}\sum_{k=0}^{m-1}
       A_{m,k}(W)_{2k}\E_a(W+2k),\label{rlc:eq:evencentral}\\
 \CC_{2m+1,a}(W;0)
 &=\frac1{(2m)!}\sum_{k=0}^{m}
       B_{m,k}(W)_{2k}\eta(W+2k,a).\label{rlc:eq:oddcentral}
\end{align}
Every summand in these formulas is entire in $W$.
\end{theorem}
\begin{proof}
For $\Rea W>0$, \eqref{rlc:eq:basefractional} at zero uses
$p_{2m}(u)/(e^u-1)$ in even depth and $p_{2m+1}(u)/(e^u+1)$
in odd depth. The ordinary Mellin integrals are
$(W)_j\zeta(W+j,a)$ or $(W)_j\eta(W+j,a)$.
In even depth $j=2k+1$, so
\[
 (W)_{2k+1}\zeta(W+2k+1,a)
     =(W)_{2k}\E_a(W+2k).
\]
This supplies the pole-free form. Entire continuation proves both
identities at all remaining orders.
\end{proof}

The first cases, useful for checking every normalization, are
\begin{align}
 \CC_{2,a}(W;0)&=\E_a(W),\label{rlc:eq:C2}\\
 \CC_{3,a}(W;0)&=\tfrac12\{(W)_2\eta(W+2,a)+\pi^2\eta(W,a)\},
 \label{rlc:eq:C3}\\
 \CC_{4,a}(W;0)&=\tfrac16\{(W)_2\E_a(W+2)+4\pi^2\E_a(W)\},
 \label{rlc:eq:C4}\\
 \CC_{5,a}(W;0)&=\tfrac1{24}\{(W)_4\eta(W+4,a)
                +10\pi^2(W)_2\eta(W+2,a)+9\pi^4\eta(W,a)\},
 \label{rlc:eq:C5}\\
 \CC_{6,a}(W;0)&=\tfrac1{120}\{(W)_4\E_a(W+4)
                +20\pi^2(W)_2\E_a(W+2)+64\pi^4\E_a(W)\}.
 \label{rlc:eq:C6}
\end{align}
An equivalent differential-operator form is
\begin{equation}\label{rlc:eq:shiftoperator}
 \CC_{r,a}(W;0)=\frac1{(r-1)!}
 \begin{cases}
 p_r(-\partial_a)\zeta(W,a),&r\text{ even},\\
 p_r(-\partial_a)\eta(W,a),&r\text{ odd}.
 \end{cases}
\end{equation}
The even polynomial has no constant term, so it annihilates the
$a$-independent zeta residue.

\subsection{Negative-order resonance is not zero substitution}
At $W=-2$, \eqref{rlc:eq:C4} gives
\begin{equation}\label{rlc:eq:resonance}
 \boxed{\displaystyle
 \CC_{4,a}(-2;0)=\frac13+\frac{2\pi^2}{3}B_2(a),
 \qquad B_2(a)=a^2-a+\frac16.}
\end{equation}
Indeed $(-2)_2\E_a(0)/6=1/3$, while
$\E_a(-2)=-2\zeta(-1,a)=B_2(a)$.
The uncombined expression $(W)_3\zeta(W+3,a)$ has a vanishing
Pochhammer factor next to a pole at $W=-2$. Setting the first factor to
zero loses the $1/3$. This is a concrete integration safeguard, not an
assertion that the inspected repository already printed that error.
